\clearpage
\section{Physical Construction of the Current Robot}
\label{sec:mechanical-design}
The robot is a differential-drive platform with two independently powered wheels, a passive rear caster and three downward-facing TCRT5000 reflectance modules. The installed computer is a \textbf{Raspberry Pi Model B Plus Rev 1.2}, with a Realtek RTL8188CUS USB Wi-Fi adapter. One Teensy 4.1 reads the sensors and generates the motor PWM. Desktop and the ASUS router are off the robot.

\subsection{Mechanical Stack and Component Placement}
The current mechanical source package is \path{cad_parts/assembly_sep24/}. Its native part files, rather than an overview snapshot, define the latest geometry. Printed parts use PETG. Figure~\ref{fig:current-physical} summarizes placement without implying measured clearances or a completed assembly test.

\begin{figure}[H]
\centering
\resizebox{.64\linewidth}{!}{\begin{tikzpicture}[font=\small,>=stealth,
 part/.style={draw,thick,rounded corners,align=center,inner sep=6pt}]
\node[part,fill=cyan!10,minimum width=8cm,minimum height=1.2cm] (upper) at (0,5.15) {Upper deck\\Teensy + clasp; power bucks; ammeter; driver mounts};
\draw[line width=5pt,black!45] (-3.3,4.5) -- (-3.3,1.8);
\draw[line width=5pt,black!45] (3.3,4.5) -- (3.3,1.8);
\node[part,fill=green!8,minimum width=5.3cm] (pi) at (0,2.75) {Pi Model B Plus + USB Wi-Fi\\Shared 2S battery};
\node[font=\bfseries,above=1pt of pi] (controller) {controller};
\node[part,fill=yellow!15,minimum width=3.8cm] (desktop) at (9.2,2.75) {Desktop};
\node[font=\bfseries,above=1pt of desktop] (operator) {operator};
\draw[dashed,thick,<->] (controller.east) -- (operator.west);
\node[part,fill=gray!18,minimum width=8cm,minimum height=.65cm] (lower) at (0,1.35) {Lower deck: motor mounts, switches and cable routing};
\draw[thick,fill=black!35] (-3.5,-.2) rectangle (-2.5,.9);
\draw[thick,fill=black!35] (2.5,-.2) rectangle (3.5,.9);
\node at (0,.4) {Two driven wheels + rear caster};
\node[part,fill=purple!8] (sensor) at (-6,1.35) {Front sensor head\\A / B / C};
\draw[thick] (sensor.east) -- (lower.west);
\draw[->,thick] (-6,.8) -- (-6,-.5);
\draw[thick] (-7,-.5) -- (4.2,-.5);
\node[font=\scriptsize,below] at (-6,-.5) {Floor / tape};
\end{tikzpicture}
}
\caption{Current mechanical arrangement, schematic and not to scale. The front sensor head reads the floor; the rear caster is passive. The removable roof is shown separated for visibility. Dimensions below are CAD values, not measurements of the assembled robot.}
\label{fig:current-physical}
\end{figure}

\begin{table}[H]
\centering\small
\caption{Current mechanical parts and their roles.}
\label{tab:current-mechanics}
\begin{tabular}{@{}L{3.2cm}L{13.0cm}@{}}
\toprule
Part & Current design description \\
\midrule
Lower deck & Nominal $110\times115\times4$\,mm, with a tapered rear. Supports the Pi, battery, motor clamps, rear caster and switches; utility openings and cable clamps provide routing access. \\
Pillar frames & Front and aft braced pairs support the upper electronics deck. The active design attaches directly to the lower deck; the old T-head joints are historical parts. \\
Upper deck & Nominal $110\times115\times3$\,mm, with Teensy rails and retaining clasp, Power A/B mounting areas, ammeter opening and inclined driver-board supports. Service openings provide access beneath components. \\
Sensor holder & Three module positions A/B/C, interpreted as left/center/right when looking forward. Mounting arms attach below the front deck. Physical channel order and working height must agree with calibration. \\
Roof & A 2\,mm plate with a 15\,mm perimeter rim and four 50\,mm columns. All four column-to-brim transitions are analytic cones. The Power A brim is 25\,mm long; the other brim arms remain 15\,mm. Nominal bracket clearance to the upper deck is 0.4\,mm. \\
\bottomrule
\end{tabular}
\end{table}

\subsection{CAD Sources and Limits of the Drawings}
\label{sec:cad-visualization}
The active native models are \path{main_lower_deck.FCStd}, \path{main_pillars.FCStd}, \path{main_upper_deck.FCStd}, \path{main_photo_sensor.FCStd} and \path{main_roof.FCStd}, all in the current package. Overview renders can lag part revisions and are not manufacturing masters. The roof's detailed dimensions are in \path{cad_parts/assembly_sep24/roof/README.md}.

Historical September 15 views and PETG estimates are retained in Appendix~\ref{app:historical-cad}; the motion model is in Appendix~\ref{sec:motion-control}. CAD checks do not establish physical fit, printed completion, assembled mass or loaded center of gravity.
